\section{Problem Formulation and Estimands}

Let $A$ denote a protected-group attribute, $Y$ a binary observed or synthetic decision, and
$S$ an optional score. For a protected group $g$ and reference group $r$, the selection rate is
$\Pr(Y=1\mid A=g)$. The adverse impact ratio used here is
\begin{equation}
  \mathrm{AIR}_{g:r}=\frac{\Pr(Y=1\mid A=g)}{\Pr(Y=1\mid A=r)}.
\end{equation}
The selection-rate gap is
\begin{equation}
  \mathrm{SRG}_{g:r}=\Pr(Y=1\mid A=g)-\Pr(Y=1\mid A=r).
\end{equation}
AIR and SRG describe association in the evaluated surface. They do not identify the cause of a gap,
the treatment effect of protected-group membership, or the legal permissibility of a decision rule.

\subsection{Internal screens}

The configured AIR point screen is \AIRThreshold{}. It is inspired by a four-fifths heuristic, but
this paper does not treat it as a fair-lending safe harbor or legal test. A point estimate and its
uncertainty interval are reported separately. The exact amplification interval is derived with
Wilson selection-rate intervals and a delta-method interval on the log ratio. Race pairwise
p-values use Holm's sequentially rejective Bonferroni adjustment; gender has one predeclared pair
\citep{wilson1927probable,oehlert1992delta,holm1979sequential}.

These intervals and p-values are conditional descriptive approximations for this generated row
surface, treating its rows as observations. They do not propagate uncertainty from fixture
construction, source-population sampling, fitted profiles, seed choice, generator version, model
selection, or deployment shift. The fixed-seed plan in Section~7 reports variation over ten
declared seeds separately; it is not a population-sampling interval.

The technical status terms mean:

\begin{description}
  \item[Within screen] the configured numerical comparison is on the expected side of its internal
    threshold;
  \item[Outside screen] the comparison triggers review under the configured technical rule;
  \item[Unavailable] required inputs or provenance are absent, so no number is reported;
  \item[Not informative] a number may be mechanically computable, but the evidence surface cannot
    support the intended interpretation; and
  \item[Not evaluated] the metric was outside the scenario's enabled capabilities.
\end{description}

\subsection{Outcome-conditioned metrics}

Equal opportunity and equalized-odds quantities require meaningful, independent ground-truth and prediction
provenance. Average odds difference (AOD) combines group differences in false-positive and
true-positive rates; equal opportunity difference (EOD) compares true-positive rates. These are
not interchangeable with selection-rate gaps \citep{hardt2016equality}.

For the exact v5.0.8 intake, every reported group has false positives and false negatives equal to
zero by construction. The resulting TPR=1 and FPR=0 surface is degenerate. FL-BSA therefore records
equalized odds as \texttt{NOT\_INFORMATIVE}; AOD and EOD are unavailable rather than reported as
zero. This distinction prevents a missing evidentiary surface from appearing to be perfect parity.

\subsection{Calibration and predictive utility}

Expected calibration error requires probability scores and outcome labels. The exact scenario has
\path{ece_enabled=false}; ECE is not evaluated. The single-run branch quality certificates also
mark predictive utility unevaluated. The robustness evidence supplies a separate amplification-only
advisory AUC diagnostic. That diagnostic is reported as a limitation because most runs are below its
configured floor; it does not convert the quality composite into a utility result
\citep{guo2017calibration,fawcett2006roc}.
